\section{Conclusions}
\label{sec:conclusion}

\phynetpy\ 1.0.0 provides a Python foundation for phylogenetic network work:
a reusable network type, an analysis layer that does not depend on any inference
method, and seven inference methods reached through two verbs over three typed
axes. Six of those methods correspond to PhyloNet analyses, and the
likelihood functions underlying two of them agree with PhyloNet's own to
$10^{-5}$ across fourteen adversarial cases.

It is a successor to PhyloNet in the sense that it continues the same line of
software and reimplements much of the same body of methods. It is not a
replacement. Two PhyloNet analyses have no counterpart here, PhyloNet command
files do not carry over, and on the one search we compared directly PhyloNet
recovered the true topology about twice as often as we did. Where those matter,
PhyloNet remains the right tool.

The case for building on \phynetpy\ is about what the next method costs. A new
criterion is a class and a registry entry, and it inherits the network
representation, the topology moves, the parameter optimiser and the diagnostics
rather than restating them.
